% cSpell:disable

\newfontfamily{\FA}[Path = fonts/]{FontAwesome.otf}

\ifnotes%
	\setbeameroption{show notes on second screen}
\fi

\title{Data-X Talk} % chktex 13

\subtitle{
	Path ORAM:\ An Extremely Simple Oblivious RAM Protocol by \emph{Emil Stefanov} and \emph{Elaine Shi} \\
	\cite{DBLP:journals/corr/abs-1202-5150} % chktex 2
}

\date{Built from \href{https://git.dbogatov.org/bu/PathORAM-Talk/commit/\version}{\emph{\version}} on \today}

\author{Dmytro Bogatov \\ dmytro@bu.edu}

\institute{Boston University}

\ifrelease%
	\logo{
		\includegraphics[
			width=1cm,
			height=1cm,
			keepaspectratio
		]{static/logo.eps}\hspace{325pt}\vspace{-30pt}
	}
\fi

\addbibresource{bibfile.bib}

\makeatletter
	\def\beamer@framenotesbegin{% at beginning of slide
		\usebeamercolor[fg]{normal text} 
			\gdef\beamer@noteitems{}%
			\gdef\beamer@notes{}% 
	}
\makeatother

\lstset{
	mathescape=true,
	numbers=left,
	tabsize=4,
	numbersep=10pt,
	breaklines=true,
	showtabs=false,
	morekeywords={for,end,do,while,if,then,else,Applies,when,return,true,false,Action},
	basicstyle=\normalfont,
	columns=fullflexible,
	frame=leftline,
	xleftmargin=15pt,
	escapechar=|
}

\hypersetup{
	pdfcopyright={Copyright (C) Dmytro Bogatov \the\year. Creative Commons Attribution-NoDerivatives 4.0 International.},
	pdflicenseurl={https://creativecommons.org/licenses/by-nd/4.0/legalcode},
	pdfauthor={Dmytro Bogatov},
	pdftitle={Data-X Talk},
	pdfsubject={Path ORAM:\ An Extremely Simple Oblivious RAM Protocol},
	pdfkeywords={Data-X, Talk, ORAM, PathORAM, Oblivious, Security, Kollios, BU, PhD, homework},
	colorlinks=true,
	linkcolor=magenta,
	urlcolor=cyan,
	citecolor=blue,
	pdfpagemode=FullScreen,
	pdfdisplaydoctitle=true,
	pdfmenubar=false,
	pdfpagelayout=SinglePage
}

\captionsetup[figure]{labelformat=empty}

\definecolor{CardinalRed}{RGB}{172,29,55} 
\definecolor{Tangerine}{RGB}{247,148,30} 

\setbeamercolor{alerted text}{fg=CardinalRed}
\setbeamercolor{example text}{fg=Tangerine}
{
	\usebeamercolor[fg]{alerted text}
	\usebeamercolor[fg]{example text}
}

\usetikzlibrary{
    shapes.geometric,
    positioning,
    fit,
    calc
}
\tikzset{
    dia/.style={
        shape=diamond,
        minimum size=2em,
    },
    dia cross/.style={
        dia,
%       path picture={
%           \draw (path picture bounding box.west) -- (path picture bounding box.east)
%                 (path picture bounding box.north) -- (path picture bounding box.south);
%       },
        append after command={
            \pgfextra%
                \draw[shorten >=\pgflinewidth,shorten <=\pgflinewidth]
                                   (\tikzlastnode.west) -- (\tikzlastnode.east); % chktex 8
                \draw[shorten >=\pgflinewidth,shorten <=\pgflinewidth]
                                 (\tikzlastnode.north) -- (\tikzlastnode.south); % chktex 8
            \endpgfextra%
        }

    }
}


